\documentclass[11pt]{article}
\usepackage[a4paper,margin=1in]{geometry}
\usepackage[T1]{fontenc}
\usepackage{lmodern}
\usepackage{amsmath}
\usepackage{microtype}
\usepackage[hidelinks]{hyperref}
\usepackage{parskip}
\usepackage{enumitem}
\setlist[itemize]{noitemsep,topsep=4pt,leftmargin=1.4em}
\pagestyle{empty}

\begin{document}

\noindent
{\large\textbf{Carl P. Zimmerman}}\\
Briar Creek Tech\\
Charlotte, North Carolina, USA

\vspace{0.6em}
\noindent\hrulefill

\vspace{0.8em}
\noindent 31 July 2026

\vspace{1.2em}
\noindent To the Editors, \textit{Physical Review D}

\vspace{0.6em}
\noindent
I submit for your consideration a short article, \textit{``The MOND acceleration scale is one half of the
dark-energy gravitational rate.''}

\vspace{0.6em}
\noindent\textbf{Context.}
The acceleration scale $a_0$ that organizes galactic rotation curves has been known since Milgrom (1983)
to lie near $cH_0$, and the coincidence has attracted a series of proposed closed forms --- a derived
$2cH_\Lambda$ from a de Sitter--Unruh argument (Milgrom 1999; also Pikhitsa 2010, Klinkhamer \& Kopp
2011), and an empirical $cH_\Lambda/2\pi$ (Milgrom 2020). This manuscript records an exact expression in
the dark-energy density with a rational coefficient, and is careful to establish what that does and does
not add to the above.

\vspace{0.6em}
\noindent\textbf{Key findings.}
\begin{itemize}
\item $a_0 = \tfrac12 c\sqrt{G\rho_\Lambda} = c^2\sqrt{\Lambda/32\pi} = cH_\Lambda/\sqrt{32\pi/3}$. The
three forms are algebraically identical, and the reduction cancels every factor of $\pi$ and the 3, which
shows that the apparent structure $\sqrt{32\pi/3}$ is only a unit conversion and carries no content.
\item On Planck 2018 inputs the expression gives $9.43\times10^{-11}\,\mathrm{m\,s^{-2}}$, which is
$0.70\%$ from the value obtained by fitting the 175 SPARC rotation curves at $\Upsilon_{[3.6]} = 0.70$ ---
a mass-to-light ratio inside the range expected for Spitzer $3.6\,\mu$m populations, so the agreement does
not require an implausible fit.
\item The relation carries a falsifiable consequence (Sec.~VI): if $a_0\propto\sqrt{\rho_\Lambda}$ then
$a_0(z)$ is exactly constant for $w=-1$ and rises to a maximum before \emph{declining} for the
DESI-preferred $w_0>-1$, $w_a<0$, in contrast to the monotonic rise predicted by $a_0\propto cH(z)$. A
measured monotonic rise would exclude the relation.
\end{itemize}

\vspace{0.4em}
\noindent\textbf{What I have not claimed, stated here because it bounds the paper.}
The interpolating function I use is \emph{identical} to Eq.~(9) of Milgrom (1999) --- not a variant of it
--- and that same paper \emph{derives} $a_0 = 2cH_\Lambda$, a factor $11.58$ larger. My contribution is
therefore a renormalization of his coefficient to match data, not a new derivation, and Sec.~II says so in
those words. Further, the coefficient here is not observationally distinguishable from Milgrom's
$cH_\Lambda/2\pi$: the two differ by $7.9\%$, which is $0.018$~dex in $g_{\rm obs}$ against a fit scatter
of $0.108$~dex. I make no claim of improved fit; the contribution is one of \emph{form}. Section~V states
that the coefficient $\tfrac12$ is fitted, not derived, and reports my own failed attempt to force it.
Section~VII hands over the outstanding tensions unprompted, including that the cluster-scale discrepancy is
$13\%$ \emph{worse} for this lower $a_0$ than for the conventional value, and that the exact interpolation
leaves a solar-system residual three orders above Earth--Mars ranging bounds.

\vspace{0.6em}
\noindent\textbf{Reproducibility and disclosure.}
Every numerical claim is produced by a public script that performs internal consistency checks and exits
with a nonzero status if any fails; the code is archived at
\url{https://doi.org/10.5281/zenodo.21724575}. Section~VIII discloses that portions of the analysis and
drafting were carried out with AI assistance: I directed the work, specified and reviewed every
load-bearing calculation, and take full responsibility including for any errors. No AI system is listed as
an author. Several intermediate claims produced during the work were found to be incorrect and were
withdrawn before submission.

\vspace{0.6em}
\noindent\textbf{Submission history.}
No version of this manuscript has previously been submitted to \textit{Physical Review D} or to any other
\textit{Physical Review} journal, nor to any other journal. It is not part of a joint submission or a
sequence of papers. A preprint is deposited at \url{https://doi.org/10.5281/zenodo.21724575}, concurrently
with this submission; it is not a journal publication. The manuscript is not posted to arXiv at the time of
submission, and I will inform the editorial office if that changes.

\vspace{0.6em}
\noindent\textbf{Recommended referees.}
I would welcome review by researchers familiar with the acceleration-scale literature, since the value of
this manuscript rests largely on whether its treatment of prior work is judged correct:
\begin{itemize}
\item Stacy S.\ McGaugh, Case Western Reserve University
\item Beno\^it Famaey, Observatoire astronomique de Strasbourg
\item Federico Lelli, INAF --- Osservatorio Astrofisico di Arcetri
\item Harry Desmond, Institute of Cosmology and Gravitation, University of Portsmouth
\end{itemize}

\vspace{0.4em}
\noindent\textbf{Excluded referees.} None.

\vspace{0.6em}
\noindent
I am an independent researcher and have no funding to declare. I am content for the manuscript to be
assessed as a short paper.

\vspace{1.2em}
\noindent With thanks for your time,

\vspace{1.4em}
\noindent Carl P. Zimmerman\\
Briar Creek Tech\\
Charlotte, North Carolina

\end{document}
